\section{Discussion}
\label{sec:discussion}

\subsection{The boundary the windows draw}

The four answers point to a single boundary that stress conditions define.
Passive provision paid in the quiet windows, where fee income is reliable and
the unhedged bands keep it. Breadth mattered more than action in the trending
windows, because a wide band carries the market exposure that a narrow band
sells away. The one high-volatility window reversed the ordering. The unhedged
rules lost heavily, while the hedged and wider positions lost less. The learned
policy sits on the defensive side of that boundary, and the premium it pays is
visible in the calm windows. On tail risk the boundary moves in its favour,
while on PnL it stays with the passive bands. Table~\ref{tab:regimes} reports
the PnL and tail loss behind these claims.

A regime label applies only to the windows assigned to it, and a single
high-volatility window cannot characterise a regime. The map therefore
describes the evaluation year rather than a fitted boundary. The study's
inference rests on the paired intervals of the primary statistics and not on
any single regime cell.

\subsection{What the hedge buys and costs}

The hedge is insurance. The study compares policies at three points along one
trade-off between protection and cost, though not as a controlled hedge
ablation. At full delta match on the opening band the tail is nearly flat and
the maximum drawdown is a fraction of the unhedged one. The bill is an average
negative hedge flow and a total PnL below the unhedged rule. At the learned
seed's partial ratios the middle of that trade-off appears. The one-shot
unhedged rules sit at the other end with a large tail and no premium, and the
two rebalancing rules trade at similar tails. Across the compared policies the
ordering is monotone in protection and in cost. Funding was positive on average
at all measured points because the short collected it, but the mark loss on a
short in a rising market and the perp fees outweighed that income.
Table~\ref{tab:aggregate} reports the corresponding PnL, Sharpe ratio and tail
loss.

The hedge does not improve the mean, it reshapes the tail, so its value depends
on the risk budget against which the LP is managed. A mandate that treats tail
risk as important as Sharpe will buy the protection that a mandate maximising
return will refuse. The pre-registered rule gives tail risk equal standing,
which is why a policy that fails the return clause and passes the tail clause
is a partial success rather than a failed idea.

\subsection{Why timing collapsed to a deployment decision}

None of the trained seeds learned to time. They deployed one band and held it.
The rebalancing benchmarks explain why. On this fee tier and in this year,
frequent recentring repeatedly realises impermanent loss, and the fees
collected do not cover it. Both active rules finish further behind the basket
than any one-shot rule, despite trading far more
(Table~\ref{tab:decomposition}). Timing is not irrelevant in principle, but it
is dominated in this market.

The study's configuration restricted the timing head after the early runs
showed the collapse, so the estimator was never asked to learn a decision that
the evidence indicates it cannot afford. That restriction is itself a finding
about the pool, and it bounds the claim. A higher fee tier, cheaper execution or
a market where price moves arrive more slowly could move the boundary, and the
width choice substitutes for timing under those conditions. The action
interface still admits the wider timing grid, so a future run can test the shift
without a redesign.

The evidence base has limits that bound the claims. The failure modes were
repaired in sequence rather than varied one at a time, so the ablation table
reports the sequence and three targeted probes, not the marginal effect of any
single design choice. The boundary is drawn from the realised 2024 windows on one pool
pair. The $0.30\,\%$ tier is measured for the rule-based ladder but not
retrained with the learned policy, and the regime with a single window is a
case study. A factorial ablation of the repairs remains open, as do a controlled
volatility and volume sweep and a learned policy on the contrast tier. All
three are re-runs of the same testbed.

\subsection{Threats to validity}

The simulator replays recorded events, so prices, liquidity and fee growth are
real rather than modelled. Its assumption that the agent is marginal is the
main simplification. The position never moves the price, and its fee share is
proportional to active liquidity. That is the standard replay assumption and it
is recorded alongside the results, but it cannot capture queueing at the tick
level or the impact of an unusually large position. The gas model charges
representative per-action units rather than transaction receipts, so gas is a
defensible floor rather than an exact quote. The hedge uses a perpetual proxy
whose funding and basis are representative rather than exact for a specific
venue.

Table~\ref{tab:threats} states the threats that bound what the study can claim
and the scope it imposes. The parity check bounds the distance between the
simulator and the backtester, not the distance between the backtester and
reality. Where the data allow, the external part of the gap is quantified by
reconciling fees and collects. Where they do not, latency, the privacy of order
flow and MEV are discussed rather than measured. We therefore claim results
about performance under risk on recorded market states under a stated cost
model. We do not claim deployable profitability, optimality of the action grid,
or that the policy would survive market conditions absent from 2022 to 2024.

\begin{table}[tbp]
\caption{Threats to validity and the scope it imposes.}
\label{tab:threats}
\centering
\small
\begin{tabular}{@{}L{0.30\linewidth}L{0.64\linewidth}@{}}
\toprule
\textbf{Threat} & \textbf{Scope} \\
\midrule
One asset pair at two fee tiers & the boundary is measured within that scope and does not automatically transfer to other pairs \\
CVaR at $5\,\%$ from few episodes & reported with bootstrap intervals and the worst window, not treated as decisive alone \\
Regime labels from fixed thresholds & a labelling device for reporting and training, not a discovered state \\
Boundary sweep over several axes & inference rests on paired intervals on the pre-registered primary metric; secondary wins are labelled as such \\
Simulator and backtester both replay recorded states & claims hold under the stated cost model, not for live execution \\
\bottomrule
\end{tabular}
\end{table}
